\setcounter{section}{19}






  \zwischenueberschrift{Restklassendarstellung für monomiale Kurven}

Es sei

\mavergleichskette
{\vergleichskette
{ M
}
{ \subseteq }{ \N
}
{  }{
}
{    }{
}
{   }{
}
}
{}{}{}
ein
\definitionsverweis {numerisches Monoid}{}{,}
das von den teilerfremden natürlichen Zahlen
\mathl{e_1  , \ldots , e_n}{} erzeugt werde. Die zugehörige Surjektion
\maabb {} { \N^n } { M \subseteq \N
} {}
führt zu einer Surjektion
\maabbeledisp {} { K[X_1  , \ldots , X_n] } { K[M]
} { X_i } { T^{e_i}} {,}
und einer abgeschlossenen Einbettung
\mathl{C=K\!-\!\operatorname{Spek}\, \left(  K[M]  \right) \hookrightarrow { {\mathbb A}_{ K }^{ n  } }}{.} Durch welche Gleichungen lässt sich $C$ beschreiben?





\inputfaktbeweis
{Affine Kurven/Monomiale Kurven/Beschreibende binomiale Gleichungen/Fakt}
{Satz}
{}
{



\faktsituation {}
\faktvoraussetzung {Es sei

\mavergleichskette
{\vergleichskette
{ M
}
{ \subseteq }{ \N
}
{  }{
}
{    }{
}
{   }{
}
}
{}{}{}
ein durch teilerfremde Elemente
\mathl{e_1 , \ldots , e_n}{} erzeugtes Untermonoid und sei
\maabb {} { \N^n } { M
} {}
die zugehörige surjektive Abbildung mit dem zugehörigen Restklassenhomomorphismus
\maabb {\varphi} { K[X_1 , \ldots , X_n] } { K[M]
} {.}}
\faktfolgerung {Dann wird das Kernideal durch

\mavergleichskettedisphandlinks
{\vergleichskettedisphandlinks
{ \ker \varphi  \!
}
{ =} {  \! \left(  \prod_{i \in I_1 } X_i^{r_i } - \prod_{i \in I_2 } X_i^{s_i } :\, I_1, I_2 \subseteq \{1, \ldots ,n\} \text{ disjunkt }, \sum_{i \in I_1 } r_ie_i = \sum_{i \in I_2 } s_ie_i  \!  \right)
}
{ } {
}
{ } {
}
{ } {
}
}
{}{}{}
\zusatzklammer {mit

\mavergleichskettek
{\vergleichskettek
{ r_i, s_i
}
{ \geq }{ 1
}
{  }{
}
{    }{
}
{   }{}
}
{}{}{}} {} {}
beschrieben.}
\faktzusatz {}
\faktzusatz {}





}
{





Dass die angegebenen Elemente zum Kernideal gehören, folgt direkt aus

\mavergleichskettedisp
{\vergleichskette
{ \varphi \left(  \prod_{i \in I_1 } X_i^{r_i }  \right)
}
{ =} { \prod_{i \in I_1 } \left(  t^{e_i }  \right)^{r_i }
}
{ =} { t^{\sum_{i \in I_1 } r_ie_i }
}
{ } {}
{ } {}
}
{}{}{.}
Für die Umkehrung sei

\mavergleichskette
{\vergleichskette
{ F
}
{ \in }{ K[X_1 , \ldots , X_n]
}
{  }{
}
{    }{
}
{   }{
}
}
{}{}{}
ein Polynom mit

\mavergleichskette
{\vergleichskette
{ \varphi (F)
}
{ = }{ 0
}
{  }{
}
{    }{
}
{   }{
}
}
{}{}{.}
Wir schreiben

\mavergleichskettedisp
{\vergleichskette
{ F
}
{ =} { \sum_\nu a_\nu X^{\nu}
}
{ } {
}
{ } {
}
{ } {
}
}
{}{}{}
\zusatzklammer {mit

\mavergleichskettek
{\vergleichskettek
{ \nu
}
{ = }{  \left(  \nu_1 , \ldots ,  \nu_n  \right)
}
{  }{
}
{    }{
}
{   }{}
}
{}{}{}} {} {.}
Daher ist

\mavergleichskettedisp
{\vergleichskette
{ \varphi(F)
}
{ =} { \sum_\nu a_\nu t^{\sum_{i=1}^n \nu_ie_i }
}
{ =} { \sum_{k=0} \left(  \sum_{ \nu : \,  \sum_{i = 1}^n \nu_i e_i = k  } a_\nu  \right)t^k
}
{ } {}
{ } {}
}
{}{}{.}
Da dieses Polynom gleich $0$ ist, müssen alle Koeffizienten $0$ sein, d.h. zu jedem $k$ gehört auch

\mavergleichskettedisp
{\vergleichskette
{ F_k
}
{ =} { \sum_{ \nu : \,  \sum_{i = 1}^n \nu_i e_i = k } a_\nu X^\nu
}
{ } {
}
{ } {
}
{ } {
}
}
{}{}{}
zum Kern. Wir können also annehmen, dass in $F$ nur Monome $X^\nu$ mit dem gleichen Wert

\mavergleichskette
{\vergleichskette
{ \sum_{i = 1}^n \nu_ie_i
}
{ = }{ k
}
{  }{
}
{    }{
}
{   }{
}
}
{}{}{}
vorkommen. Betrachten wir ein solches Monom aus $F$, sagen wir $X^\nu$
\zusatzklammer {mit

\mavergleichskettek
{\vergleichskettek
{ a_\nu
}
{ \neq }{ 0
}
{  }{
}
{    }{
}
{   }{
}
}
{}{}{}} {} {.}
Es muss in $F$ mindestens noch ein weiteres Monom, sagen wir $X^\mu$, vorkommen, da ein einzelnes Monom nicht auf $0$ abgebildet wird. Wir schreiben

\mavergleichskettedisp
{\vergleichskette
{ F
}
{ =} { a_\nu  \left(  X^\nu -X^\mu  \right)  + \left(  F- a_\nu X^\nu + a_\nu X^\mu  \right)
}
{ } {
}
{ } {
}
{ } {
}
}
{}{}{.}
Im Summanden rechts kommt $X^\nu$ nicht mehr vor, und es kommt auch kein neues Monom hinzu. In
\mathl{X^\nu-X^\mu}{} können wir diejenigen Variablen, die beidseitig auftreten, so weit ausklammern, dass sich ein Ausdruck der Form

\mavergleichskettedisp
{\vergleichskette
{ X^\nu-X^\mu
}
{ =} { X_1^{b_1 } \cdots X_n^{b_n } \left(  \prod_{i \in I_1 } X_i^{r_i }- \prod_{i \in I_2 } X_i^{s_i }  \right)
}
{ } {
}
{ } {
}
{ } {
}
}
{}{}{}
mit disjunkten $I_1$ und $I_2$ und mit

\mavergleichskette
{\vergleichskette
{ \sum_{i \in I_1 }e_ir_i
}
{ = }{ \sum_{i \in I_2 }e_is_i
}
{  }{
}
{    }{
}
{   }{
}
}
{}{}{}
ergibt. Der linke Summand in obiger Beschreibung von $F$ gehört also zu dem von den angegebenen Binomen erzeugten Ideal und wir können mit dem rechten Summanden, in dem ein Monom weniger vorkommt, fortfahren.




}



Die im vorstehenden Satz auftretenden Gleichungen nennt man \stichwort {binomiale Gleichungen} {.} Die einfachsten binomialen Gleichungen sind von der Bauart
\zusatzklammer {
\mavergleichskettek
{\vergleichskettek
{ i
}
{ \neq }{ j
}
{  }{
}
{    }{
}
{   }{
}
}
{}{}{}} {} {}

\mavergleichskettedisp
{\vergleichskette
{ X_i^{e_j/\operatorname{ggT}(e_i,e_j)}
}
{ =} { X_j^{e_i/\operatorname{ggT}(e_i,e_j)}
}
{ } {
}
{ } {
}
{ } {
}
}
{}{}{.}
Im Fall von ebenen monomialen Kurven ist das auch die einzige Gleichung.



\inputfaktbeweis
{Ebene affine Kurven/Monomiale Kurve/Beschreibende Gleichung/Fakt}
{Korollar}
{}
{



\faktsituation {}
\faktvoraussetzung {Es sei $C$ die durch
\mathl{t \mapsto (t^{e_1 }, t^{e_2 })=(x,y)}{}
\zusatzklammer {mit $e_1,e_2$ teilerfremd} {} {}
gegebene monomiale ebene Kurve.}
\faktfolgerung {Dann ist

\mavergleichskettedisp
{\vergleichskette
{C
}
{ =} { V(X^{e_2 }-Y^{e_1 })
}
{ } {
}
{ } {
}
{ } {
}
}
{}{}{.}}
\faktzusatz {}
\faktzusatz {}





}
{





Dies folgt sofort aus

Satz 19.1
.




}



Bei monomialen Raumkurven lassen sich die beschreibenden Gleichungen auch noch einigermaßen einfach bestimmen, da man immer eine Variable isolieren kann.

\inputbeispiel{}
{





$\,$

\bild{ \begin{center}
\includegraphics[width=5.5cm]{ \bildeinlesungpng {Twisted_cubic_curve} {png} }
\end{center}
\bildtext {} }
\bildlizenz { Twisted cubic curve.png } {Claudio Rocchini} {} {Commons} {CC-BY-SA-3.0} {}



Es sei

\mavergleichskette
{\vergleichskette
{ C
}
{ \subset }{ { {\mathbb A}_{  K  }^{  3   } }
}
{  }{
}
{    }{
}
{   }{
}
}
{}{}{}
die \anfuehrung{gedrehte Kubik}{,} also das Bild der monomialen Abbildung, die durch
\mathl{t \mapsto \left(  t,t^2,t^3  \right)}{} gegeben ist. Diese Kurve ist isomorph zu einer affinen Geraden und insbesondere glatt. Das beschreibende Ideal ist nach

Satz 19.1

gleich

\mavergleichskettedisp
{\vergleichskette
{ {\mathfrak a}
}
{ =} { \left(  Y-X^2,Z-X^3,Y^3-Z^2,Z-XY  \right)
}
{ =} { \left(  Y-X^2,Z-X^3  \right)
}
{ } {
}
{ } {
}
}
{}{}{.}
Die beiden letzten Idealerzeuger sind dabei überflüssig, da sie sich durch die beiden anderen ausdrücken lassen. Insgesamt ist also

\mavergleichskettedisp
{\vergleichskette
{ C
}
{ =} { V \left(  Y-X^2,Z-X^3  \right)
}
{ } {
}
{ } {
}
{ } {
}
}
{}{}{.}
Die Bilder von $C$ unter den drei verschiedenen Projektionen sind

\mathdisp {C_1 = V \left(  Z^2-Y^3  \right) ,\, C_2 = V \left(  Z-X^3  \right) , \, C_3 = V \left(  Y-X^2  \right)} { . }

Dabei sind $C_2$ und $C_3$ isomorph zur affinen Geraden
\zusatzklammer {als Graph einer Abbildung} {} {,}
während $C_1$ die singuläre Neilsche Parabel ist.





}

\inputbeispiel{}
{





Es sei $C$ die durch

\mathdisp {t \longmapsto (t^3,t^4,t^5) = (x,y,z)} {  }

gegebene monomiale Kurve. Für jede der drei Variablen müssen wir gemäß

Satz 19.1

schauen, welche Potenzen davon, wenn man die $t$-Potenzen substituiert, sich auch als Monom in den beiden anderen Variablen ausdrücken lassen.

Zunächst haben wir die Gleichungen, in denen jeweils nur zwei Variablen vorkommen. Das sind

\mathdisp {Y^3 = X^4,\, Z^3 = X^5, \, Z^4 = Y^5} { . }

Hier kann es, wie im ebenen Fall, immer nur eine Beziehung geben.

In den Relationen, wo alle drei Variablen beteiligt sind, kommt eine der Variablen allein vor. Starten wir mit $X$. Zunächst lassen sich $X$ und $X^2$ nicht durch die anderen Variablen ausdrücken, dafür haben wir

\mavergleichskette
{\vergleichskette
{ X^3
}
{ = }{ T^9
}
{ = }{ YZ
}
{    }{
}
{   }{
}
}
{}{}{.}
Eine andere
\zusatzklammer {davon unabhängige} {} {}
Kombination ist nicht möglich. Grundsätzlich impliziert eine mehrfache Darstellung

\mavergleichskette
{\vergleichskette
{ X^k
}
{ = }{ Y^{i}Z^{j}
}
{ = }{ Y^{a}Z^{b}
}
{    }{
}
{   }{
}
}
{}{}{,}
dass man zwischen Potenzen von $Y$ und von $Z$ eine Beziehung hat, da man ja die kleineren Potenzen rauskürzen kann. Da wir alle Relationen mit nur zwei Variablen schon aufgelistet haben, liefert eine Potenz von $X$ immer nur maximal eine neue Relation. Wir behaupten, dass wir für $X$ alleinstehend schon fertig sind. Ist nämlich

\mavergleichskette
{\vergleichskette
{ X^k
}
{ = }{ Y^{i}Z^{j}
}
{  }{
}
{    }{
}
{   }{
}
}
{}{}{,}
so ist

\mavergleichskette
{\vergleichskette
{ k
}
{ \geq }{ 3
}
{  }{
}
{    }{
}
{   }{
}
}
{}{}{.}
Bei

\mavergleichskette
{\vergleichskette
{ i
}
{ = }{ 0
}
{  }{
}
{    }{
}
{   }{
}
}
{}{}{}
oder

\mavergleichskette
{\vergleichskette
{ j
}
{ = }{ 0
}
{  }{
}
{    }{
}
{   }{
}
}
{}{}{}
haben wir die Gleichungen schon aufgelistet. Es sei also

\mavergleichskette
{\vergleichskette
{ i,j
}
{ \geq }{ 1
}
{  }{
}
{    }{
}
{   }{
}
}
{}{}{.}
Dann kann man aber mittels der Gleichung

\mavergleichskette
{\vergleichskette
{ X^3
}
{ = }{ YZ
}
{  }{
}
{    }{
}
{   }{
}
}
{}{}{}
die Exponenten in der Gleichung kleiner machen
\zusatzklammer {indem man den Exponenten von $X$ um $3$ reduziert und die Exponenten von $Y$ und von $Z$ um $1$} {} {.}

Für $Y$ hat man sofort die Gleichung

\mavergleichskette
{\vergleichskette
{ Y^2
}
{ = }{ ZX
}
{  }{
}
{    }{
}
{   }{
}
}
{}{}{,}
mit der man wieder alle anderen Gleichungen reduzieren kann.

Für $Z$ hat man

\mavergleichskette
{\vergleichskette
{ Z^2
}
{ = }{ X^2Y
}
{  }{
}
{    }{
}
{   }{
}
}
{}{}{}
und

\mavergleichskette
{\vergleichskette
{ Z^3
}
{ = }{ XY^3
}
{  }{
}
{    }{
}
{   }{
}
}
{}{}{.}
Es gibt keine kleineren Monome in $X$ und $Y$, die man als Potenz von $Z$ ausdrücken kann. Daher kann man jede andere Relation mittels einer von diesen auf eine frühere zurückführen.

Insgesamt haben wir also für die Kurve $C$ die Gleichungen

\mavergleichskettedisp
{\vergleichskette
{ C
}
{ =} { V( Y^3-X^4, Z^3-X^5, Z^4-Y^5,
 \mathdisplaybruch X^3-YZ, Y^2-XZ, Z^2-X^2Y ,Z^3-XY^3 )
}
{ } {
}
{ } {
}
{ } {
}
}
{}{}{.}





}





  \zwischenueberschrift{Ganzheit}


\inputdefinition
{}
{





Es seien
\mathkor {} {R} {und} {S} {}
\definitionsverweis {kommutative Ringe}{}{}
und sei

\mavergleichskette
{\vergleichskette
{R
}
{ \subseteq }{ S
}
{  }{
}
{    }{
}
{   }{
}
}
{}{}{}
eine Ring\-erweiterung. Für ein Element

\mavergleichskette
{\vergleichskette
{ x
}
{ \in }{ S
}
{  }{
}
{    }{
}
{   }{
}
}
{}{}{}
heißt eine Gleichung der Form

\mavergleichskettedisp
{\vergleichskette
{ x^n+ r_{n-1}x^{n-1} + r_{n-2}x^{n-2} + \cdots + r_1x +r_0
}
{ =} { 0
}
{ } {
}
{ } {
}
{ } {
}
}
{}{}{,}
wobei die Koeffizienten \mathind { r_{i}  } {  i=0 , \ldots , n-1   }{,} zu $R$ gehören, eine \definitionswort {Ganzheitsgleichung}{} für $x$.





}


\inputdefinition
{}
{





Es seien
\mathkor {} {R} {und} {S} {}
\definitionsverweis {kommutative Ringe}{}{}
und

\mavergleichskette
{\vergleichskette
{R
}
{ \subseteq }{ S
}
{  }{
}
{    }{
}
{   }{
}
}
{}{}{}
eine Ring\-erweiterung. Ein Element

\mavergleichskette
{\vergleichskette
{ x
}
{ \in }{ S
}
{  }{
}
{    }{
}
{   }{
}
}
{}{}{}
heißt \definitionswort {ganz}{}
\zusatzklammer {über $R$} {} {,}
wenn $x$ eine
\definitionsverweis {Ganzheitsgleichung}{}{}
mit Koeffizienten aus $R$ erfüllt.





}


\inputdefinition
{}
{





Es seien
\mathkor {} {R} {und} {S} {}
\definitionsverweis {kommutative Ringe}{}{}
und

\mavergleichskette
{\vergleichskette
{R
}
{ \subseteq }{ S
}
{  }{
}
{    }{
}
{   }{
}
}
{}{}{}
eine Ring\-erweiterung. Dann nennt man die Menge der Elemente

\mavergleichskette
{\vergleichskette
{ x
}
{ \in }{ S
}
{  }{
}
{    }{
}
{   }{
}
}
{}{}{,}
die
\definitionsverweis {ganz}{}{}
über $R$ sind, den \definitionswort {ganzen Abschluss}{} von $R$ in $S$.





}


\inputdefinition
{}
{





Es seien
\mathkor {} {R} {und} {S} {}
\definitionsverweis {kommutative Ringe}{}{}
und sei

\mavergleichskette
{\vergleichskette
{R
}
{ \subseteq }{ S
}
{  }{
}
{    }{
}
{   }{
}
}
{}{}{}
eine Ring\-erweiterung. Dann heißt $S$ \definitionswort {ganz}{} über $R$, wenn jedes Element

\mavergleichskette
{\vergleichskette
{ x
}
{ \in }{ S
}
{  }{
}
{    }{
}
{   }{
}
}
{}{}{}
\definitionsverweis {ganz}{}{}
über $R$ ist.





}

$S$ ist genau dann ganz über $R$, wenn der ganze Abschluss von $R$ in $S$ gleich $S$ ist.





\inputfaktbeweis
{Kommutative Ringtheorie/Ganzheit/Ganzes Element/Charakterisierung/Fakt}
{Lemma}
{}
{



\faktsituation {}
\faktvoraussetzung {Es seien $R$ und $S$
\definitionsverweis {kommutative Ringe}{}{}
und

\mavergleichskette
{\vergleichskette
{R
}
{ \subseteq }{ S
}
{  }{
}
{    }{
}
{   }{
}
}
{}{}{}
eine Ringerweiterung.}
\faktuebergang {Für ein Element

\mavergleichskette
{\vergleichskette
{ x
}
{ \in }{ S
}
{  }{
}
{    }{
}
{   }{
}
}
{}{}{}
sind folgende Aussagen äquivalent.}
\faktfolgerung {\aufzaehlungdrei{ $x$ ist
\definitionsverweis {ganz}{}{}
über $R$.
}{Es gibt eine $R$-Unteralgebra $T$ von $S$ mit

\mavergleichskette
{\vergleichskette
{ x
}
{ \in }{ T
}
{  }{
}
{    }{
}
{   }{
}
}
{}{}{}
und die ein endlicher $R$-Modul ist.
}{Es gibt einen endlichen $R$-Untermodul $M$ von $S$, der einen
\definitionsverweis {Nichtnullteiler}{}{}
aus $S$ enthält, mit

\mavergleichskette
{\vergleichskette
{ xM
}
{ \subseteq }{ M
}
{  }{
}
{    }{
}
{   }{
}
}
{}{}{.}}}
\faktzusatz {}
\faktzusatz {}





}
{





(1) $\Rightarrow$ (2). Wir betrachten die von den Potenzen von $x$ erzeugte $R$-Unteralgebra
\mathl{R[x]}{} von $S$, die aus allen polynomialen Ausdrücken in $x$ mit Koeffizienten aus $R$ besteht. Aus einer Ganzheitsgleichung

\mavergleichskettedisp
{\vergleichskette
{ x^n+ r_{n-1}x^{n-1} + r_{n-2}x^{n-2} + \cdots +  r_1x +r_0
}
{ =} { 0
}
{ } {
}
{ } {
}
{ } {
}
}
{}{}{}
ergibt sich

\mavergleichskettedisp
{\vergleichskette
{ x^n
}
{ =} { - r_{n-1}x^{n-1} - r_{n-2}x^{n-2} - \cdots -  r_1x -r_0
}
{ } {
}
{ } {
}
{ } {
}
}
{}{}{.}
Man kann also $x^n$ durch einen polynomialen Ausdruck von einem kleineren Grad ausdrücken. Durch Multiplikation dieser letzten Gleichung mit $x^{i}$ kann man jede Potenz von $x$ mit einem Exponenten $\geq n$ durch einen polynomialen Ausdruck von einem kleineren Grad ersetzen. Insgesamt kann man dann aber all diese Potenzen durch polynomiale Ausdrücke vom Grad
\mathl{\leq n-1}{} ersetzen. Damit ist

\mavergleichskettedisp
{\vergleichskette
{ R[x]
}
{ =} { R + Rx + Rx^2 + \cdots +  Rx^{n-2} + Rx^{n-1}
}
{ } {
}
{ } {
}
{ } {
}
}
{}{}{}
und die Potenzen
\mathl{x^0=1,x^1,x^2 , \ldots , x^{n-1}}{} bilden ein endliches $R$-Modul-Erzeugendensystem von

\mavergleichskette
{\vergleichskette
{T
}
{ = }{ R[x]
}
{  }{
}
{    }{
}
{   }{
}
}
{}{}{.}

(2) $\Rightarrow$ (3). Es sei

\mavergleichskette
{\vergleichskette
{ x
}
{ \in }{ T
}
{ \subseteq }{ S
}
{    }{
}
{   }{
}
}
{}{}{,}
$T$ eine $R$-Unteralgebra, die als $R$-Modul endlich erzeugt sei. Dann ist

\mavergleichskette
{\vergleichskette
{ xT
}
{ \subseteq }{ T
}
{  }{
}
{    }{
}
{   }{
}
}
{}{}{,}
und $T$ enthält den Nichtnullteiler $1$.

(3) $\Rightarrow$ (1). Sei

\mavergleichskette
{\vergleichskette
{M
}
{ \subseteq }{ S
}
{  }{
}
{    }{
}
{   }{
}
}
{}{}{}
ein endlich erzeugter $R$-Untermodul mit

\mavergleichskette
{\vergleichskette
{ xM
}
{ \subseteq }{ M
}
{  }{
}
{    }{
}
{   }{
}
}
{}{}{.}
Es seien
\mathl{y_1 , \ldots , y_n}{} erzeugende Elemente von $M$. Dann ist insbesondere
\mathl{xy_i}{} für jedes $i$ eine $R$-Linearkombination der \mathind {  y_j   } { j=1 , \ldots , n   }{.} Dies bedeutet

\mavergleichskettedisp
{\vergleichskette
{ x y_i
}
{ =} { \sum_{j = 1}^n r_{ij} y_j
}
{ } {
}
{ } {
}
{ } {
}
}
{}{}{}
mit

\mavergleichskette
{\vergleichskette
{ r_{ij}
}
{ \in }{ R
}
{  }{
}
{    }{
}
{   }{
}
}
{}{}{,}
oder, als Matrix geschrieben,

\mavergleichskettedisp
{\vergleichskette
{ x \begin{pmatrix}  y_1  \\ y_2 \\ \vdots\\ y_n   \end{pmatrix}
}
{ =} { \begin{pmatrix} r_{1,1}  &  r_{1,2}  &  \ldots  &  r_{1,n}  \\  r_{2,1}  &  r_{2,2}  &  \ldots  &  r_{2,n}  \\  \vdots  & \vdots  &  \ddots  &  \vdots   \\  r_{n,1}   &  r_{n,2}  &  \ldots  &  r_{n,n} \end{pmatrix} \begin{pmatrix}  y_1  \\ y_2 \\ \vdots\\ y_n   \end{pmatrix}
}
{ } {
}
{ } {
}
{ } {
}
}
{}{}{.}
Dies schreiben wir als

\mavergleichskettedisp
{\vergleichskette
{ 0
}
{ =} { \begin{pmatrix}  x-r_{1,1}  &  - r_{1,2}  &  \ldots  &  - r_{1,n}  \\  - r_{2,1}  &  x -r_{2,2}  &  \ldots  &  - r_{2,n}  \\  \vdots  & \vdots  &  \ddots  &  \vdots   \\  -r_{n,1}   &  - r_{n,2}  &  \ldots  &  x- r_{n,n} \end{pmatrix} \begin{pmatrix}  y_1  \\ y_2 \\ \vdots\\ y_n   \end{pmatrix}
}
{ } {
}
{ } {
}
{ } {
}
}
{}{}{.}
Nennen wir diese Matrix $A$
\zusatzklammer {die Einträge sind aus $S$} {} {,}
und sei
\mathl{A ^{ \operatorname{adj} }}{} die
\definitionsverweis {adjungierte Matrix}{}{.}
Dann gilt

\mavergleichskette
{\vergleichskette
{ A ^{ \operatorname{adj} } A y
}
{ = }{ 0
}
{  }{
}
{    }{
}
{   }{
}
}
{}{}{}
\zusatzklammer {$y$ bezeichne den Vektor \mathlk{(y_1 , \ldots , y_n)}{}} {} {}
und
nach der
Cramerschen Regel

ist

\mavergleichskette
{\vergleichskette
{ A ^{ \operatorname{adj} } A
}
{ = }{ (\det  A )E_n
}
{  }{
}
{    }{
}
{   }{
}
}
{}{}{,}
also gilt

\mavergleichskettedisp
{\vergleichskette
{ ((\det  A ) E_n) y
}
{ =} { 0
}
{ } {
}
{ } {
}
{ } {
}
}
{}{}{.}
Es ist also

\mavergleichskette
{\vergleichskette
{ (\det  A ) y_j
}
{ = }{ 0
}
{  }{
}
{    }{
}
{   }{
}
}
{}{}{}
für alle $j$ und damit

\mavergleichskettedisp
{\vergleichskette
{ (\det  A ) z
}
{ =} { 0
}
{ } {
}
{ } {
}
{ } {
}
}
{}{}{}
für alle

\mavergleichskette
{\vergleichskette
{z
}
{ \in }{ M
}
{  }{
}
{    }{
}
{   }{
}
}
{}{}{.}
Da $M$ nach Voraussetzung einen Nichtnullteiler enthält, muss

\mavergleichskette
{\vergleichskette
{ \det  A
}
{ = }{ 0
}
{  }{
}
{    }{
}
{   }{
}
}
{}{}{}
sein. Die Determinante ist aber ein normierter polynomialer Ausdruck in $x$ vom Grad $n$, sodass eine Ganzheitsgleichung vorliegt.




}






\inputfaktbeweis
{Kommutative Ringtheorie/Ganzheit/Ganzer Abschluss/Ring/Fakt}
{Korollar}
{}
{



\faktsituation {}
\faktvoraussetzung {Es seien
\mathkor {} {R} {und} {S} {}
\definitionsverweis {kommutative Ringe}{}{}
und sei

\mavergleichskette
{\vergleichskette
{R
}
{ \subseteq }{ S
}
{  }{
}
{    }{
}
{   }{
}
}
{}{}{}
eine
\definitionsverweis {Ring\-erweiterung}{}{.}}
\faktfolgerung {Dann ist der
\definitionsverweis {ganze Abschluss}{}{}
von $R$ in $S$ eine $R$-Unteralgebra von $S$.}
\faktzusatz {}
\faktzusatz {}





}
{





Die Ganzheitsgleichungen \mathind { X-r  } { r \in R   }{,} zeigen, dass jedes Element aus $R$ ganz über $R$ ist. Es seien

\mavergleichskette
{\vergleichskette
{ x_1
}
{ \in }{ S
}
{  }{
}
{    }{
}
{   }{
}
}
{}{}{}
und

\mavergleichskette
{\vergleichskette
{ x_2
}
{ \in }{ S
}
{  }{
}
{    }{
}
{   }{
}
}
{}{}{}
ganz über $R$. Nach
der
Charakterisierung der Ganzheit

gibt es endliche $R$-Unteralgebren

\mavergleichskette
{\vergleichskette
{ T_1,T_2
}
{ \subseteq }{ S
}
{  }{
}
{    }{
}
{   }{
}
}
{}{}{}
mit

\mavergleichskette
{\vergleichskette
{ x_1
}
{ \in }{ T_1
}
{  }{
}
{    }{
}
{   }{
}
}
{}{}{}
und

\mavergleichskette
{\vergleichskette
{ x_2
}
{ \in }{ T_2
}
{  }{
}
{    }{
}
{   }{
}
}
{}{}{.}
Es sei
\mathl{y_1 , \ldots , y_n}{} ein $R$-Erzeugendensystem von $T_1$ und
\mathl{z_1 , \ldots , z_m}{} ein $R$-Erzeugendensystem von $T_2$. Wir können annehmen, dass

\mavergleichskette
{\vergleichskette
{ y_1
}
{ = }{ z_1
}
{ = }{ 1
}
{    }{
}
{   }{
}
}
{}{}{}
ist. Betrachte den endlich erzeugten $R$-Modul

\mavergleichskettedisp
{\vergleichskette
{ T
}
{ =} { T_1 \cdot T_2
}
{ =} { \langle y_iz_j, \, i= 1 , \ldots ,  n, \, j = 1 , \ldots , m \rangle
}
{ } {
}
{ } {
}
}
{}{}{,}
der offensichtlich
\mathl{x_1+x_2}{} und
\mathl{x_1x_2}{}
\zusatzklammer {und $1$} {} {}
enthält. Dieser $R$-Modul $T$ ist auch wieder eine $R$-Algebra, da für zwei beliebige Elemente gilt

\mavergleichskettedisp
{\vergleichskette
{ \left(  \sum r_{ij} y_iz_j  \right) \left(  \sum s_{kl} y_kz_l  \right)
}
{ =} { \sum  r_{ij}s_{kl} y_iy_k z_jz_l
}
{ } {
}
{ } {
}
{ } {
}
}
{}{}{,}
und für die Produkte gilt

\mavergleichskette
{\vergleichskette
{ y_iy_k
}
{ \in }{ T_1
}
{  }{
}
{    }{
}
{   }{
}
}
{}{}{}
und

\mavergleichskette
{\vergleichskette
{ z_jz_l
}
{ \in }{ T_2
}
{  }{
}
{    }{
}
{   }{
}
}
{}{}{,}
sodass diese Linearkombination zu $T$ gehört. Dies zeigt, dass die Summe und das Produkt von zwei ganzen Elementen wieder ganz ist. Deshalb ist der ganze Abschluss ein Unterring von $S$, der $R$ enthält. Also liegt eine $R$-Unteralgebra vor.




}





\inputdefinition
{}
{





Es seien $R$ und $S$
\definitionsverweis {kommutative Ringe}{}{}
und

\mavergleichskette
{\vergleichskette
{R
}
{ \subseteq }{ S
}
{  }{
}
{    }{
}
{   }{
}
}
{}{}{}
eine
\definitionsverweis {Ringerweiterung}{}{.}
Man nennt $R$ \definitionswort {ganz-abgeschlossen}{} in $S$, wenn der
\definitionsverweis {ganze Abschluss}{}{}
von $R$ in $S$ gleich $R$ ist.





}
